\chapter{Installing and starting}\label{ch:install}

\section{The packages}

Everything the program needs is inside each package: the tracker drivers and the other files it
loads at run time are in the \texttt{resources} folder next to the program. (The \texttt{.deb} of
the Linux architectures without an AppImage asks the package manager for the libraries instead of
carrying them, section~\ref{sec:linux-arch}.)

\begin{center}
\begin{tabular}{@{}>{\raggedright\arraybackslash}p{3.0cm}>{\raggedright\arraybackslash}p{4.6cm}>{\raggedright\arraybackslash}p{6.2cm}@{}}
\toprule
System & Package & How to start it \\ \midrule
Linux x86\_64 (Debian 11, Ubuntu 20.04 and newer) & \texttt{Ritmo-Linux-\allowbreak x86\_64.AppImage} &
  make it executable (\texttt{chmod +x}) and run it \\
Windows 64 bit & \texttt{Ritmo-Windows-x64-Setup.exe}, \texttt{Ritmo-Windows-x64.zip} &
  run the installer (it asks whether to install for all the users or for the current one, and offers a desktop icon and the
  \texttt{.rmt} file type), or unzip the archive anywhere and run \texttt{Ritmo.exe} (the program is not signed: Windows may ask to confirm) \\
macOS 12 or later, Apple Silicon or Intel & \texttt{Ritmo-macOS-arm64.dmg}, \texttt{Ritmo-macOS-x86\_64.dmg} &
  drag \texttt{Ritmo.app} to Applications; the first time open it with the right button and
  \menu{Open} (the program is not signed) \\
Linux aarch64: Raspberry Pi 3/4/5 with a 64 bit system (Raspberry Pi OS 13, Debian 13, Ubuntu 24.04 and newer) &
  \texttt{Ritmo-Linux-\allowbreak aarch64.AppImage} & as the x86\_64 one \\
Linux riscv64 (Debian 13 and newer) & \texttt{ritmo\_}\emph{version}\texttt{\_riscv64.deb}, \texttt{Ritmo-Linux-\allowbreak riscv64.tar.gz} &
  install the \texttt{.deb} with \texttt{apt}, or unpack the archive (section~\ref{sec:linux-arch}) \\
Linux ppc64: PowerPC 64 bit big endian (Debian ports, \emph{sid}) & \texttt{ritmo\_}\emph{version}\texttt{\_ppc64.deb}, \texttt{Ritmo-Linux-\allowbreak ppc64.tar.gz} &
  as riscv64 \\
\bottomrule
\end{tabular}
\end{center}

\RITMO{} can also be built from the source with CMake and Qt (appendix~\ref{app:building}).

\begin{note}
The packages of the versions up to 2.2.1 were called \texttt{RMT-*} and the program \texttt{rmt}
(\texttt{Rmt.exe}). From the next releases the packages are \texttt{Ritmo-*} and the program
\texttt{ritmo} (\texttt{Ritmo.exe}).
\end{note}

\section{Linux, architecture by architecture}\label{sec:linux-arch}

\texttt{uname -m} tells the architecture of the machine: \texttt{x86\_64}, \texttt{aarch64}, \texttt{riscv64} or \texttt{ppc64}.

\subsection{x86\_64 and aarch64 (Raspberry Pi): the AppImage}

An AppImage is one file with the program and everything it needs (Qt, PortAudio, RtMidi). It only needs the
\texttt{libfuse2} library to start (on Debian 13 and Ubuntu 24.04 the package is called \texttt{libfuse2t64}):
\begin{lstlisting}
sudo apt install libfuse2                  # libfuse2t64 on Debian 13 and Ubuntu 24.04
chmod +x Ritmo-Linux-x86_64.AppImage       # Raspberry Pi: Ritmo-Linux-aarch64.AppImage
./Ritmo-Linux-x86_64.AppImage mysong.rmt
\end{lstlisting}
Without FUSE the file can still run, unpacking itself in a temporary folder:
\begin{lstlisting}
./Ritmo-Linux-x86_64.AppImage --appimage-extract-and-run mysong.rmt
\end{lstlisting}
The x86\_64 AppImage runs on Debian 11, Ubuntu 20.04 and newer. The aarch64 one needs glibc 2.38 or newer: Raspberry Pi OS 13
(\emph{trixie}), Debian 13, Ubuntu 24.04 and newer, on a Raspberry Pi 3, 4 or 5 with a 64 bit system (\texttt{uname -m} must say
\texttt{aarch64}: on a 32 bit system it says \texttt{armv7l} and there is no package).

\subsection{riscv64 and ppc64: the .deb}

For Debian and the systems derived from it, riscv64 needs Debian 13 or newer, and ppc64 (PowerPC 64 bit \emph{big endian}) the
\emph{sid} release of Debian ports, the only one that has this architecture. The \texttt{.deb} contains the program and
\texttt{apt} installs the libraries it needs (Qt~6, PortAudio, RtMidi...) by itself:
\begin{lstlisting}
sudo apt install ./ritmo_2.5_riscv64.deb   # ppc64: ./ritmo_2.5_ppc64.deb
ritmo mysong.rmt                           # also in the menu of the desktop
\end{lstlisting}
The \texttt{./} before the name is needed: without it \texttt{apt} looks for a package of that name in the repositories.
The program goes in \texttt{/usr/lib/ritmo}, with the link \texttt{/usr/bin/ritmo}. To remove it:
\begin{lstlisting}
sudo apt remove ritmo
\end{lstlisting}

\subsection{riscv64 and ppc64: the archive}

For the other distributions, or where \texttt{apt} is not used, the \texttt{.tar.gz} carries Qt~6, PortAudio and RtMidi (folders
\texttt{lib} and \texttt{plugins}); of the machine it only uses glibc and the graphics (OpenGL) and sound (ALSA) libraries of its
drivers. The \texttt{README.txt} inside lists the Debian packages for these:
\begin{lstlisting}
tar xzf Ritmo-Linux-riscv64.tar.gz         # ppc64: Ritmo-Linux-ppc64.tar.gz
cd Ritmo-Linux-riscv64
cat README.txt                             # the packages the machine needs
sudo apt install <the packages of README.txt>
./ritmo mysong.rmt
\end{lstlisting}
\texttt{ritmo} is a script that sets where the libraries and the plug-ins are: it can be started from any folder, and the
\texttt{resources} folder must stay next to it. Nothing is installed: to remove it, delete the folder. The glibc of the machine
cannot be older than the one the archive was built with (Debian 13 for riscv64, \emph{sid} for ppc64).

\begin{note}
There are no packages for PowerPC 32 bit and for ppc64le (little endian PowerPC). The packages of riscv64 and ppc64 have been built
and started (without a display) in an emulator, not on real machines; the engine of the ppc64 build gives the same results as the
x86\_64 one on all the songs of the distribution. If you try them, please tell the result in the issues of the project.
\end{note}

\section{Windows and macOS in detail}\label{sec:win-mac}

\subsection{Windows 64 bit: the installer}

Run \texttt{Ritmo-Windows-x64-Setup.exe}. The program is not signed, so Windows (SmartScreen) may ask to confirm: choose
\menu{More info} and then \menu{Run anyway}. The installer asks whether to install for all the users (it needs the
administrator) or only for the current one, and offers a desktop icon and the association of the \texttt{.rmt} files with
\RITMO{}, so that a double click on a song opens it. It adds an entry to the Start menu and to the list of the programs of Windows,
from which it is removed. The settings are kept in the registry of the user, so a new installation does not lose them.

\subsection{Windows 64 bit: the .zip}

\texttt{Ritmo-Windows-x64.zip} has the same files without the installer: unzip it in any folder (also on a USB stick) and run
\texttt{Ritmo.exe}; the \texttt{resources} folder and the DLLs of Qt must stay next to it. Nothing is installed and no file type is
associated: to remove it, delete the folder (the settings stay in the registry, \texttt{HKEY\_CURRENT\_USER\textbackslash Software\textbackslash ritmo-atari.org}).

\subsection{macOS 12 or later: the disk image}

Open \texttt{Ritmo-macOS-arm64.dmg} (Apple Silicon) or \texttt{Ritmo-macOS-x86\_64.dmg} (Intel; it also runs on Apple Silicon
through Rosetta~2) and drag \texttt{Ritmo.app} to the \texttt{Applications} folder shown in the window. The application is signed
only ad hoc, not notarized: the first time, open it with the right button (or Control and click) and \menu{Open}, then confirm.
After that it starts normally. It carries Qt and the libraries, nothing else needs to be installed. To remove it, move
\texttt{Ritmo.app} to the bin; the settings are in \texttt{\textasciitilde/Library/Preferences/org.ritmo-atari.ritmo.plist}.

\section{Starting the program}

Start \RITMO{} like any other program. The first start opens the window with an empty song.
The program accepts a song on the command line, which is opened at the start:
\begin{lstlisting}
ritmo mysong.rmt
\end{lstlisting}

With the option \texttt{/SCRIPT:} the program runs a script without showing its window
(chapter~\ref{ch:scripting}):
\begin{lstlisting}
ritmo /SCRIPT:build.rmtscript
\end{lstlisting}

With \texttt{--scale=}\emph{n} (also \texttt{/scale:}\emph{n}), from 100 to 300, the program starts with that interface size
(section~\ref{sec:window}) for this session only: the settings are not changed.
\begin{lstlisting}
ritmo --scale=200 mysong.rmt
\end{lstlisting}

\section{The window}\label{sec:window}

The window is made for a laptop with a 1366$\times$768 screen: that is its size the first time the program starts (reduced to what the screen has, if
the screen is smaller), and the whole screen of a mono song fits in it, with the registers of the POKEY next to the song area (\menu{View $\rightarrow$ Pokey Chip Registers}, on by
default). A stereo song also fits, with its eight tracks and the song area, but the registers of the POKEY, which need about 650 more pixels on
the right, are left out: to see them, make the window wider (at least about 1720 pixels) or use a larger screen. After that the program remembers the size and the position of the window
and opens it the same way next time (if the screen is smaller than before, the window is reduced to fit).

The \emph{interface size} (section~\ref{sec:options}) is chosen the first time too, from the size of the screen: 100\,\% on the screens up to Full HD,
200\,\% on a screen of 2560$\times$1440, 300\,\% on a 4K screen with no scaling of the system. Only whole multiples are chosen, so that the pixels of the
tracker stay square. The choice can be changed in the \menu{Options} dialog, and the settings taken over from \RMT{} keep the size they had.

\section{Sound}

The sound goes out through PortAudio, to the default output device of the system. If the output cannot be opened
the program prints a message and goes on without sound, so that songs can still be edited. If the sound
crackles or is distorted, choose a larger \emph{audio buffer} in the configuration (section~\ref{sec:options}).
The \texttt{RMT\_AUDIO\_BUFFER\_MS} environment variable (the size of the buffer in milliseconds, up to 50) sets
another size for one session without changing the settings.

\section{Where the settings are kept}

The configuration and the tuning belong to the user, not to the program folder, so that an update
or a new installation keeps them:

\begin{center}
\begin{tabular}{@{}ll@{}}
\toprule
Linux & \texttt{\textasciitilde/.config/ritmo-atari.org/ritmo.conf} (\texttt{\$XDG\_CONFIG\_HOME}) \\
Windows & the registry, \texttt{HKEY\_CURRENT\_USER\textbackslash Software\textbackslash ritmo-atari.org\textbackslash ritmo} \\
macOS & \texttt{\textasciitilde/Library/Preferences/org.ritmo-atari.ritmo.plist} \\
\bottomrule
\end{tabular}
\end{center}

At the first start, if there are no settings of \RITMO{} yet, the settings of \RMT{} are taken over: those
saved by it in the same places with the name \texttt{rmt}, or an \texttt{rmt.ini} and a
\texttt{tuning.ini} next to the program.
